\documentclass[11pt]{article}
\usepackage[letterpaper,margin=1in]{geometry}
\usepackage[T1]{fontenc}
\usepackage{lmodern,microtype}
\usepackage{amsmath,amssymb,amsthm,mathrsfs}
\usepackage{xcolor}
\usepackage[colorlinks=true,linkcolor=blue!40!black,citecolor=blue!40!black,urlcolor=blue!40!black]{hyperref}
\usepackage{fancyhdr}
\pagestyle{fancy}
\fancyhf{}
\fancyhead[L]{\small First-order twistor integrability}
\fancyhead[R]{\small Working draft}
\fancyfoot[C]{\thepage}
\setlength{\headheight}{14pt}
\setlength{\emergencystretch}{2em}
\numberwithin{equation}{section}
\newtheorem{proposition}{Proposition}[section]
\DeclareMathOperator{\Ric}{Ric}
\DeclareMathOperator{\tr}{tr}
\newcommand{\CP}{\mathbb{CP}}

\title{First-order twistor integrability and the\\
smooth four-dimensional Poincar\'e conjecture}
\author{Edward B. Baker III}
\date{Working draft --- 25 November 2026}
% Prepared for Edward B. Baker III with OpenAI Codex assistance,
% 23 November 2026 (America/New_York). Selected model and effort unavailable.
% The inherited opening and companion input are pinned in ../DRAFT_HISTORY.md.

\begin{document}
\maketitle
\thispagestyle{plain}

% Revision from f178fbb against research through 5ff7014. Serving settings
% of the coordinating session unavailable. Companion checks selected Sol/ultra.
\begin{abstract}
Three local sections of the positive twistor sphere give pointwise probes
of chiral Weyl curvature. We derive an integrated estimate from their
Frobenius tensors and identify the corresponding curvature component for a
positive spinor. On a degenerating homotopy four-sphere we track where
such estimates must hold to imply a strict energy improvement on the
original smooth manifold. Marked regions carrying physical volume can now
be found in the subcritical band. Extending their markings to a separating
collar or a nonround closed base remains the geometric question.
\end{abstract}

\section{Introduction}\label{sec:introduction}

Isothermal coordinates provide a simple model for using geometry to
recognize the smooth structure of a sphere. Every smooth Riemannian metric
on an oriented real surface can be written locally as
\[
 g=e^{2u}(dx^2+dy^2)=e^{2u}|dz|^2,\qquad z=x+iy.
\]
Orientation-preserving transitions between these charts are holomorphic,
so the metric determines a Riemann surface. If the underlying smooth
surface is closed and homotopy equivalent to \(S^2\), it is orientable
and simply connected. Uniformization then identifies this compact Riemann
surface biholomorphically with \(\CP^1\); the resulting map is a
diffeomorphism to the standard smooth two-sphere~\cite{Donaldson2011}.
This gives a geometric proof of the two-dimensional smooth Poincar\'e
statement. Isothermal coordinates supply the local conformal description,
while uniformization supplies the global identification; the existence of
the local charts alone would not establish the conclusion.

In three dimensions, the analogous statement is that every closed smooth
three-manifold homotopy equivalent to \(S^3\) is diffeomorphic to the
standard three-sphere. Hamilton's approach replaces conformal coordinates
by a geometric heat equation, the Ricci flow,
\[
 \partial_t g=-2\Ric(g).
\]
After fixing the diffeomorphism freedom, this is a nonlinear parabolic
evolution of the metric. Hamilton proved that in dimension three the
volume-normalized flow starting from strictly positive Ricci curvature
converges to a metric of constant positive sectional
curvature~\cite{Hamilton1982}. General initial metrics can develop
singularities, so the broader program required both their analysis and
a way to continue the evolution. Perelman's entropy and noncollapsing
estimates supplied decisive control of singularity models~\cite{Perelman2002}.
His construction of Ricci flow with surgery removes appropriate neck
regions and replaces them by controlled caps~\cite{Perelman2003Surgery}.
For a homotopy three-sphere, finite-time extinction together with the
topology of these surgeries yields the smooth identification with
\(S^3\)~\cite{Perelman2003Extinction}. The lesson for higher dimensions is
that a local smoothing equation must be accompanied by global control of
degeneration and its topological consequences.

In higher dimensions, Smale's handle-cancellation methods, underlying
$h$-cobordism theory, show that every closed smooth homotopy $n$-sphere
with $n\geq5$ is homeomorphic to $S^n$~\cite{Smale1961}.
Milnor's exotic seven-spheres show that the same topological sphere can
admit distinct smooth structures~\cite{Milnor1956}.
Dimension four lies outside the range of these smooth handle-cancellation
arguments. Freedman proved that every closed topological homotopy
four-sphere is homeomorphic to \(S^4\)~\cite{Freedman1982}. The smooth
four-dimensional Poincar\'e conjecture asks whether every closed smooth
homotopy four-sphere is diffeomorphic to the standard \(S^4\). This remains
open. Major approaches include handle calculus and surgery descriptions
that seek to simplify homotopy
spheres, and knot-theoretic obstructions that could distinguish a candidate
exotic sphere~\cite{FreedmanGompfMorrisonWalker2010}.
Geometric approaches seek a distinguished metric or connection
whose existence would force the standard smooth structure. The present
work pursues this direction, with anti-self-dual (ASD) conformal geometry
as a possible target. Fine~\cite{Fine2011} introduced a gradient flow for
definite \(SO(3)\) connections aimed at producing ASD Einstein metrics of
nonzero scalar curvature. This offers a concrete
four-dimensional evolution framework; its application to the smooth
Poincar\'e problem would still require suitable initial connections on
arbitrary homotopy spheres and global control of the flow.

The special geometry of four dimensions explains the appeal of these
variables. Unlike a surface metric, a general four-dimensional metric has
Weyl curvature that cannot be removed by conformal rescaling or a change
of coordinates. At the same time, on an oriented Riemannian four-manifold
the Hodge star satisfies \(*^2=1\) on real two-forms, producing two real
rank-three bundles,
\[
 \Lambda^2T^*M=\Lambda_g^+\oplus\Lambda_g^-.
\]
Atiyah, Hitchin, and Singer~\cite{AHS1978} used this chiral splitting to
connect conformal geometry and gauge fields. In our orientation convention,
their twistor integrability criterion says that \(W^+=0\) precisely when
the natural almost complex
structure on the sphere bundle \(S(\Lambda_g^+)\) is integrable.
The resulting complex threefold encodes a real four-dimensional conformal
geometry through a family of projective lines. Penrose's nonlinear-graviton
construction supplies the corresponding
holomorphic picture and, with additional data, its Ricci-flat
specialization~\cite{Penrose1976}. These are four-dimensional analogues of
the relation between conformal geometry and complex analysis on surfaces,
with a curvature integrability condition in place of an automatic local
complex structure.

Gauge theory provides a complementary reason to work with connections.
Donaldson invariants extract information about smooth four-manifolds from
moduli spaces of anti-self-dual gauge connections~\cite{DonaldsonKronheimer1990}.
Their success motivates the use of gauge fields in smooth topology,
although it does not settle the
homotopy-sphere case, where the second cohomology vanishes. Twistor methods
also reach beyond the self-dual sector. Witten~\cite{Witten1978} describes
the full classical Yang--Mills equations using holomorphic bundles on the
space of complex null lines, with specified triviality and third-order
formal extension conditions. This is a local
complex flat-spacetime correspondence, not an existence theorem for curved
Riemannian four-manifolds. Together these ideas motivate generalized
coordinate descriptions in which maps depend on both a point and a chiral
direction, and curvature is carried by connection data. For a flat-space
precursor based on generalized conformal maps acting on Yang--Mills fields,
see~\cite{Baker2023} and references therein.

The starting question here is whether the conformal structure can be
constructed directly by solving for integrable local sections of
\(S(\Lambda_g^+)\). The surface argument suggests what such a
construction would have to replace. If \(u\) is harmonic with
\(du\ne0\) on an oriented surface, \(d(*du)=0\) and locally one can
choose \(v\) with \(dv=*du\). The two differentials are independent,
orthogonal, and equal in norm, giving isothermal coordinates.
In four real dimensions \(*du\) has degree three. The closer analogue
is the Beltrami formulation, in which the coordinate functions are
annihilated by a complex first-order distribution. The involutivity
condition on the real metric twistor space is now a curvature
condition for the base metric.

Let \(M\) be a closed, connected, smooth, oriented homotopy four-sphere.
Throughout we use positive-definite Riemannian signature. For a smooth
metric \(g\), write \(W_g^\pm\) for the Weyl curvature operators on
\(\Lambda_g^\pm\). In the operator Hilbert--Schmidt convention,
\begin{equation}\label{eq:chiral-weyl-energy}
 E_\pm(g)=\int_M |W_g^\pm|_{\mathrm{op,HS}}^2\,dV_g,\qquad
 E_+(g)-E_-(g)=12\pi^2\tau(M).
\end{equation}
The signature of \(M\) is zero, so \(W^+=0\) implies \(W^-=0\) by
smoothness and \eqref{eq:chiral-weyl-energy}. In that case the local
conformal charts develop into \(S^4\). Simple connectivity gives a
global developing map. Compactness makes this local diffeomorphism
a covering, which is a diffeomorphism since \(S^4\) is simply
connected~\cite[\S4, eq.~(16)]{Hirzebruch1971Signature}
\cite[pp.~47--50]{SchoenYau1988Conformal}. Constructing a conformal
class with \(W^+=0\) on an arbitrary such \(M\) is therefore equivalent
to the smooth classification question.

We develop three formulations of the proposed construction. The first
uses several projective spinor directions to express involutivity on
each patch. The second uses one positive-chiral twistor spinor as a
sufficient curvature condition. The third uses the compatibility
equation between a chiral connection and its metric two-forms. These
equations are first order in the respective independent fields. The
condition \(W^+=0\) in the metric alone is second order. The modified
Bach flow differentiates curvature twice and is fourth order in the
metric~\cite[eqs.~(1.1)--(1.2)]{BahuaudGuentherIsenbergMazzeo2025Bach}.
The proposed direct route must produce a single conformal class
compatible with the local solutions. Section~\ref{sec:three-defects}
quantifies the fixed-metric obstruction using three sections or one
spinor. Section~\ref{sec:marked-accounting} follows that obstruction
through marked truncations of a sequence and identifies the remaining
geometric input.

\section{Local twistor sections and the Weyl quartic}
\label{sec:local-sections}

For this section allow any smooth oriented Riemannian four-manifold.
The twistor fiber of \(Z_g=S(\Lambda_g^+)\) parametrizes positive
orthogonal almost complex structures \(J\) on a tangent space. In
the two-form norm a structure \(J\) corresponds to the unit
self-dual form \(\omega_J/\sqrt{2}\), with
\(\omega_J(X,Y)=g(JX,Y)\). Let \(\mathcal J\) denote the natural
almost complex structure on \(Z_g\). It acts vertically by the
complex structure of the sphere and horizontally by the \(J\)
represented by that point, using the Levi--Civita splitting. In this
chirality, the Atiyah--Hitchin--Singer integrability criterion is
\(W^+=0\)~\cite[Theorem~4.1, pp.~435--436]{AHS1978}.

Locally, in auxiliary complex coordinates \(z^j\) for a smooth
coordinate patch on the real six-manifold \(Z_g\), the \((0,1)\)
distribution of \(\mathcal J\) can be written as the graph of a
matrix \(\mu\) after adapting the coordinates at a point. It is
spanned by
\begin{equation}\label{eq:beltrami}
 L_j=\partial_{\bar z^j}+\sum_{a=1}^3\mu_j^{\,a}\partial_{z^a},
 \qquad j=1,2,3.
\end{equation}
One seeks three complex coordinate functions \(f^a\) satisfying
\(L_jf^a=0\) with six real independent differentials. The
compatibility condition is closure of the span of the \(L_j\) under
brackets. In \eqref{eq:beltrami}, the bracket has components only
along the auxiliary \(\partial_{z^a}\), so closure requires its
coefficients to vanish. For the surface Beltrami equation a complex
line distribution is automatically involutive. On \(Z_g\), closure
is equivalent to \(W^+=0\) by the stated twistor criterion.

A section of \(Z_g\) over an open set of the base is called
integrable here when its associated \(J\) is integrable on that open
set. A single integrable section can occur with nonzero \(W^+\).
For example, the Fubini--Study metric is K\"ahler for the complex
structure of \(\CP^2\), with its complex orientation. That complex
structure gives an integrable section, while
\(\tau(\CP^2)=1\) and \eqref{eq:chiral-weyl-energy} imply
\(E_+>0\)~\cite[eq.~(16)]{Hirzebruch1971Signature}
\cite[Theorem~B]{ChangGurskyYang2003}.

\begin{proposition}[Three local directions]
\label{prop:three-directions}
Let \(U\) be an open set in a smooth oriented Riemannian
four-manifold. The curvature \(W^+\) vanishes on \(U\) if and only
if every point of \(U\) has a neighborhood in \(U\) with three
integrable sections \(J_1,J_2,J_3\) of \(Z_g\) for the same metric,
satisfying \(J_i\ne\pm J_j\) pointwise for \(i\ne j\).
\end{proposition}

\begin{proof}
For an integrable \(J\), put \(\mathcal D=T^{0,1}_J\). This is an
involutive complex two-plane distribution, totally null for the
complex-bilinear extension of \(g\). For \(X,Y,Z\) sections of
\(\mathcal D\), the terms in the Koszul formula for
\(2g(\nabla_XY,Z)\) vanish because \(\mathcal D\) is null and
involutive. Since the plane is maximal null,
\(\nabla_XY\) is again a section of \(\mathcal D\). Thus
\(g(R(X,Y)X,Y)=0\). The Schouten terms in the curvature
decomposition vanish on these null arguments, so \(W(X,Y,X,Y)=0\).

Where \(X,Y\) span \(\mathcal D\), their wedge determines a null
line in \(\Lambda^+_{\mathbb C}\), identifying vectors with
covectors by \(g\). As the positive twistor direction varies, these
lines are the null conic in
\(\mathbb P(\Lambda^+_{\mathbb C})\). In a local spin frame, a
representative \(\xi(\pi)\) of the line is quadratic in the
components of the projective spinor parameter \([\pi]\).
Consequently the pairing
\[
 q_{W^+}(\pi)=\langle W^+\xi(\pi),\xi(\pi)\rangle_{\mathbb C}
\]
is a homogeneous quartic, where the pairing is complex-bilinear.
The vanishing on an involutive plane is one root of this
quartic~\cite[Remark~5.2, Proposition~5.3 and eqs.~(14)--(20)]{GoverHillNurowski2011GoldbergSachs}.
The six directions \(J_i,-J_i\) are distinct because the three
sections are distinct modulo sign. They are all integrable, so
\(q_{W^+}\) vanishes identically. To see that it detects the entire
block, evaluate the same quadratic pairing on the null vectors
\(u_i+\mathrm{i}u_j\) and \(u_i-\mathrm{i}u_j\) in an orthonormal
real basis \((u_1,u_2,u_3)\) of \(\Lambda^+\). The off-diagonal
entries of \(W^+\) vanish and its diagonal entries are equal.
Trace-freeness gives \(W^+=0\).

Conversely, suppose \(W^+=0\) on \(U\). Then \(Z_g|_U\) is complex.
Through any chosen point in a twistor fiber, take a local complex
surface with tangent complementary to the vertical complex line.
Projection to the base is a local diffeomorphism and yields a
section. At all points of the section the projection intertwines
the induced complex structure of the surface with the represented
\(J\), by the horizontal and vertical description of
\(\mathcal J\)~\cite[Theorem~4.1 and its remarks]{AHS1978}.
This \(J\) is integrable. Choose three directions distinct up to
sign at the base point and shrink the common domain to retain
that separation.
\end{proof}

\section{Projective spinors and twistor spinors}
\label{sec:spinors}

\subsection{Projective spinors and Frobenius}

On a spin coordinate patch, write \(T_{\mathbb C}M\simeq
S^-\otimes S^+\), with \(S^+\) chosen so that real self-dual
two-forms act on it. For a nonzero positive spinor field \(\pi\),
let
\[
 \mathcal D_\pi=S^-\otimes\operatorname{span}_{\mathbb C}(\pi).
\]
This plane field represents one positive twistor direction.
The null planes are complex subspaces of the complexified tangent
bundle of a real Riemannian manifold. This construction takes
place on the real base and uses its Riemannian spin connection.
The line \([\pi]\) defines an integrable section if and only if
\begin{equation}\label{eq:projective-frobenius}
 (\nabla^+_X\pi)\wedge\pi=0
 \quad\text{for all }X\in\mathcal D_\pi.
\end{equation}
Indeed, the Koszul argument in Proposition~\ref{prop:three-directions}
says that involutivity makes \(\nabla_X\mathcal D_\pi\) lie in
\(\mathcal D_\pi\) for \(X\in\mathcal D_\pi\). By the tensor product
connection this is exactly \eqref{eq:projective-frobenius}.
Conversely the same preservation, together with zero torsion,
implies involutivity. The derivative of a nonvanishing rescaling of
\(\pi\) contributes only a term in the line, so the equation is
projective. A local spin lift always suffices for this
description~\cite[\S3 and \S7.1]{GoverHillNurowski2011GoldbergSachs}.

For a fixed metric the equations test the available directions.
The proposed construction instead seeks a common unknown conformal
metric and three lines satisfying \eqref{eq:projective-frobenius}
on each patch. The local criterion identifies the target exactly.
It remains to find a method for solving the coupled equations
with compatible conformal data on overlaps.

\subsection{The chiral twistor equation}

The ordinary Dirac operator \(D_g\) enters the twistor operator by
removing the Clifford trace of the covariant derivative. In the
convention \(X\cdot X=-|X|_g^2\), a positive-chiral twistor spinor
satisfies
\begin{equation}\label{eq:twistor-spinor}
 P_X\psi_+ :=\nabla_X\psi_++\tfrac14 X\cdot D_g\psi_+=0
 \quad\text{for all real tangent vectors }X.
\end{equation}
Equation~\eqref{eq:twistor-spinor} is first order. Its curvature
compatibility is \(W(X,Y)\cdot\psi_+=0\) for real \(X,Y\).
Friedrich obtains this from a connection for which
\((\psi_+,D_g\psi_+)\) is parallel~\cite[eq.~(1.1) and pp.~61--62]{Friedrich1990TwistorSpinors}.

Only the self-dual part of the real two-form \(W(X,Y)\) acts on
\(S^+\). A nonzero real element of \(\Lambda^+\) acts as a nonzero
trace-free skew-Hermitian endomorphism of a complex two-dimensional
space. Its eigenvalues are opposite nonzero imaginary scalars, so the
action is invertible. At any point with \(\psi_+\ne0\), the
self-dual part of \(W(X,Y)\) therefore vanishes for every \(X,Y\).
Projections of the decomposable real two-forms span \(\Lambda^+\),
which gives \(W^+=0\) at the point. On a connected spin manifold a
nontrivial twistor spinor has dense nonzero locus. If it vanishes on
an open set, the parallel pair vanishes there, hence everywhere.
Thus a nontrivial global positive-chiral twistor spinor gives
\(W^+=0\) on the entire connected base.

A homotopy four-sphere is spin since
\(H^2(M;\mathbb Z/2)=0\). The construction would need to allow zeros
of \(\psi_+\), as considered in Section~\ref{sec:overlap}.
Existence of \(\psi_+\) with \eqref{eq:twistor-spinor} imposes the
prolongation compatibility conditions in addition to testing the
Weyl block. The global harmonic-spinor equation \(D_g\psi=0\)
has zero kernel already on the round \(S^4\), by Lichnerowicz
vanishing for positive scalar curvature~\cite[p.~460]{AHS1978}.

\section{Chiral connection and metric two-forms}
\label{sec:connection}

Let \(\Sigma=(\Sigma_1,\Sigma_2,\Sigma_3)\) be a smooth triple of
two-forms in the round-compatible frame convention and let \(\nu\)
be a positive volume form satisfying
\begin{equation}\label{eq:frame}
 \Sigma_i\wedge\Sigma_j=2\delta_{ij}\nu.
\end{equation}
The span of the two-forms is a positive three-plane for the wedge
pairing, so it determines the conformal metric with this span as
\(\Lambda_g^+\). The volume choice \(dV_g=\nu\) fixes the
scale~\cite[Definitions~1.3--1.6]{Fine2011}. Under oriented
orthogonal frame changes, \(\Sigma\) is a two-form valued in an
oriented Euclidean rank-three bundle \(E\).
Let \(A\) be an independent \(SO(3)\) connection on \(E\), and use
the companion manuscript's curvature and bracket
conventions~\cite[Sections~1.3 and~1.6]{Fine2011}
\cite[adapted frame and compatibility in Section~4]{Baker2026ChiralMetricEvolution}.
Define \(C\) by
\begin{equation}\label{eq:positive-block}
 C_{ij}\,dV_g=-\tfrac12 F_A^i\wedge\Sigma_j .
\end{equation}
The negative sign agrees with \(F_{\Gamma^+}=-K\Sigma\) in constant
sectional curvature \(K\). The proposed first-order system is
\begin{equation}\label{eq:chiral-first-order}
 d_A\Sigma=0,\qquad C-\tfrac13(\tr C)I=0.
\end{equation}

The first equation identifies \(A\) with the chiral Levi--Civita
connection \(\Gamma^+\) of \(g\). To see this, write
\(b=A-\Gamma^+\). Since \(d_{\Gamma^+}\Sigma=0\), the difference is
\(d_A\Sigma=[b,\Sigma]\). In a normalized positive frame, the
pointwise map \(\mathscr C b=*[b,\Sigma]\) on \(E\)-valued
one-forms satisfies \(\mathscr C^2+\mathscr C-2I=0\).
It is invertible by the companion's compatibility
calculation~\cite[Section~4.4, compatibility proposition]{Baker2026ChiralMetricEvolution}.
With \(A=\Gamma^+\), the chiral curvature decomposition in these
conventions gives
\[
 F_A=-\bigl(W^++\tfrac{s_g}{12}I\bigr)\Sigma
      -B_{\mathrm{Ric}}\Sigma^-,
\]
where the coefficient \(B_{\mathrm{Ric}}\) is the trace-free Ricci
block and \(\Sigma^-\) is an anti-self-dual frame with the same
norms. Wedge product with \(\Sigma_j\) kills the latter
term. Thus \(C=W^++s_g I/12\), and the second equation of
\eqref{eq:chiral-first-order} is precisely
\(W^+=0\)~\cite[Section~5, chiral curvature decomposition]{Baker2026ChiralMetricEvolution}.
Conversely, a smooth metric with \(W^+=0\) supplies the normalized
triple locally, and \(A=\Gamma^+\) satisfies both equations.

The unknowns in \eqref{eq:chiral-first-order} include the metric
two-forms and the connection, so each term uses at most one
derivative of an independent field. Eliminating \(A\) recovers the
second-order metric equation. Fine's first-order perfect-connection
equation \(Q=I\) uses a definite curvature triple and gives an ASD
Einstein metric of nonzero scalar curvature in its stated
setting~\cite[Sections~1.3 and~1.6]{Fine2011}. The positivity
condition in \eqref{eq:frame} is on the metric two-forms. The
trace and mixed Ricci block in \eqref{eq:chiral-first-order}
keep their values from the compatible metric.

The covariant closure equation also explains a difference from
\(dv=*du\). One has \(d_A^2=[F_A,\cdot]\), and ordinary local
exactness for \(d\) is replaced by a compatibility problem for
\((\Sigma,A)\). On a homotopy four-sphere a globally defined
closed self-dual two-form is exact since \(H^2(M;\mathbb R)=0\).
Its squared norm integrates to zero by Stokes' theorem, so it
vanishes. The frame-valued unknown \(\Sigma\) allows the local
three-planes to be described on a bundle with nontrivial
transitions.

\section{Conformal overlaps and the conditioning of three directions}
\label{sec:overlap}

An integrable section need only be local. In fact, a single global
section of \(Z_g\) would give an almost complex structure on
the homotopy four-sphere. Its Chern classes would satisfy
\[
 \langle c_1^2,[M]\rangle
 =\langle p_1(TM)+2e(TM),[M]\rangle
 =3\tau(M)+2\chi(M)=4.
\]
The first Chern class vanishes because \(H^2(M;\mathbb Z)=0\),
a contradiction. A construction by twistor spinors must likewise
allow zeros since a nowhere-zero positive spinor would select a
global projective direction. The local sections can be chosen
independently on different patches. The corresponding conformal
metrics must agree across overlaps, with spin and frame data
compared through their transition laws. Proposition
\ref{prop:three-directions} then gives \(W^+=0\) throughout the
resulting conformal class, to which the recognition argument in
Section~\ref{sec:introduction} applies.

A first test would use two overlapping patches of the round
\(S^4\). Choose three separated sections on each patch and study
the linearization of the coupled metric-and-section equations.
Specify the allowable diffeomorphism, conformal scale, and spin
frame changes, and ask which infinitesimal conformal discrepancies
on the overlap can be removed. The chiral system
\eqref{eq:chiral-first-order} provides a second description to
compare with the metric linearization. At a round metric of
sectional curvature \(K>0\), the companion calculation gives the
\(E_+\) Hessian on transverse-traceless variations
\(\tfrac18(L+2K)(L+4K)\), where \(L=\nabla^*\nabla\).
This fixes a benchmark for the metric part of an overlap
analysis~\cite[round metric Hessian factorization]{Baker2026ChiralMetricEvolution}
\cite[Theorem~3.10]{GurskyViaclovsky2015}. The choice of an
overlap estimate or construction remains open.

There is a related quantitative question. At a point, take three
twistor directions with separation at least \(\varepsilon>0\)
modulo the antipodal map, for the round sphere metric on the fiber.
Let \(\xi_i\) be representatives of their three complex null
lines in \(\Lambda^+_{\mathbb C}\), each of Hermitian norm one.
For a real symmetric trace-free block \(W^+\) there is a uniform
algebraic bound
\begin{equation}\label{eq:quartic-estimate}
 |W^+|_{\mathrm{op,HS}}\ \le\
 C_\varepsilon
 \left(\sum_{i=1}^3
 \bigl|\langle W^+\xi_i,\xi_i\rangle_{\mathbb C}\bigr|^2\right)^{1/2}.
\end{equation}
Indeed, the three complex evaluations determine the block, because
their vanishing supplies the three conjugate evaluations as well.
The quartic root argument then gives injectivity. The minimum of
the norm of the evaluation map on unit blocks is positive
uniformly over the compact set of separated triples. The magnitudes
of the evaluations are independent of the choices of unit phase.

\section{Measuring the defect of three sections}\label{sec:three-defects}

The estimate takes a particularly simple form when the three sections
are orthogonal in the twistor sphere. Write their axes as an orthonormal
basis \(n_1,n_2,n_3\) of \(\Lambda^+\), and set
\(\xi_i=(n_j+\mathrm i n_k)/\sqrt2\) for cyclic \(i,j,k\).
For a symmetric trace-free matrix \(W=(w_{ij})\) in this basis, write
\(d_j=w_{jj}\) and \(q_i=|\langle W\xi_i,\xi_i\rangle_{\mathbb C}|\).
The complex bilinear evaluation is \((d_j-d_k)/2+\mathrm i w_{jk}\).
Since \(\sum d_i=0\), summing its squared modulus gives
\begin{equation}\label{eq:orthogonal-samples}
 \sum_i q_i^2=\tfrac34\sum_i d_i^2+\sum_{j<k}w_{jk}^2
 \ \ge\ \tfrac12|W|_{\rm HS}^2 .
\end{equation}
Equality holds for any block with zero diagonal. Thus the reconstruction
constant is uniform even where the eigenlines of \(W\) cannot be chosen
smoothly.

For a positive orthogonal almost complex structure \(J\), put
\(P=(1+\mathrm iJ)/2,\ Q=1-P,\ F(X,Y)=Q[PX,PY]\). Torsion freeness
makes this bracket tensorial. Give the complexified tangent its
Hermitian norm, counting each exterior pair once, and use the complex
bilinear metric for null pairings. If \(X,Y\) form a unitary basis of
\(\mathcal D=\mathrm{im}\,P\), the Schouten part vanishes on the null
plane. In the connection preserving the splitting into \(\mathcal D\)
and its conjugate, the off-diagonal part of the Levi--Civita connection is
\(K=(\nabla P)(P-Q)\). It has norm \(|\nabla J|/2\). The null Koszul
formula expresses the off-diagonal connection on \(\mathcal D\) through
\(F\). Differentiating it gives the curvature on the unit null bivector
\(X\wedge Y\) as
\[
 g_{\mathbb C}((\nabla_X F)(X,Y),Y)
 -g_{\mathbb C}((\nabla_Y F)(X,Y),X)
 -g_{\mathbb C}(K_{F(X,Y)}Y,X).
\]
The last term accounts for the torsion of the splitting connection.
The first two terms together have magnitude at most
\(\sqrt2\,|\nabla F|\). Cauchy--Schwarz in the two quantities
\(|\nabla F|\) and \(|\nabla J||F|\) now bounds the squared evaluation
by \(\tfrac94(|\nabla F|^2+|\nabla J|^2|F|^2)\)
\cite{Baker2026TwistorSections}.

With \(W=W_g^+\), multiply \eqref{eq:orthogonal-samples} and the
last bound by a subordinate partition of unity \(\{\rho_a\}\) and
integrate. Orthogonal triples \(J_{ai}\) can be chosen independently
on the open sets \(U_a\) of a finite cover, giving
\[
 E_+(g)\le\tfrac92\sum_{a,i}\int_{U_a}\rho_a
  (|\nabla F_{ai}|^2+|\nabla J_{ai}|^2|F_{ai}|^2)\,dV_g .
\]
There is no differentiation of the partition in this argument. On a
smooth closed oriented homotopy four-sphere in a positive Yamabe class,
a right hand side strictly below \(4\pi^2\) gives smooth recognition.
Indeed the signature identity yields equality of the two chiral
integrals, each in Hilbert--Schmidt norm, and summing ordered
four-index components gives four times their sum. The full tensor
integral then satisfies the strict sphere theorem
\cite[Theorem~A and Remark~2]{ChangGurskyYang2003}.

The scale of the lower bound can also be tested from above. Let
\(I(g)\) be the infimum of this weighted integral of the defect
density (without the prefactor \(9/2\)) over finite local covers,
partitions and orthogonal triples. At a point \(p\) choose a
conformal representative agreeing to first order with \(g\), whose
Schouten tensor vanishes at \(p\), by prescribing the Hessian of the
logarithmic conformal factor there. Radial transport of the axes in
this representative makes their first derivatives vanish at \(p\).
The derivative of its connection matrix in a radial frame is one half
of curvature at the center. Only the positive curvature block acts on
the axes; it is \(W^+\) there. With the exterior-pair norm this gives
\(\sum_i |\nabla F_i|^2(p)\le4|W^+|_{\rm HS}^2(p)\).
By continuity the local inequality holds to arbitrary additive error
in small patches about each center. Choosing a finite subcover and
a partition, then sending the error to zero, gives
\[
 \tfrac29 E_+(g)\le I(g)\le4E_+(g).
\]
The same shrinking-patch argument gives zero for the infimum with
density \(\sum_i|F_i|^4\), since all the brackets vanish at their
chosen centers \cite{Baker2026TwistorSections}. This explains the
derivative in the functional even on one smooth metric.

\subsection{A spinor supplies the axes}

Let \(j\) be the parallel antiunitary quaternionic structure of
\(S^+\). Where a smooth positive-chirality spinor \(\psi\) is nonzero,
take \(\pi_1=\psi,\ \pi_2=(\psi+j\psi)/\sqrt2,
\pi_3=(\psi+\mathrm i j\psi)/\sqrt2\). For unit representatives the
axis inner product is \(2|\langle u,v\rangle|^2-1\). The three
pairwise squared inner products are \(1/2\), giving the axes required
above. For a single line write \(u=\pi/|\pi|\), and let \(q\) project
onto its orthogonal complement. On
\(\mathcal D=S^-\otimes\mathbb C u\), the Clifford trace correction in
\((P_{\rm tw}\pi)_X=\nabla_X\pi+c(X)D\pi/4\) takes values in
\(\mathbb C u\). For \(X=a\otimes u,Y=b\otimes u\), torsion freeness
therefore gives
\[
 F(X,Y)= b\otimes v_X-a\otimes v_Y,\qquad
 v_X=q(P_{\rm tw}\pi)_X/|\pi|.
\]
Choosing \(a,b\) unitary makes the output terms orthogonal. The three
spinors have equal norms, and \(P_{\rm tw}\) commutes with \(j\) on
real tangent vectors. Consequently
\(\sum_i |F_i|^2\le 3|P_{\rm tw}\psi|^2/|\psi|^2\)
\cite{Baker2026TwistorSections}. This controls the undifferentiated
brackets. Curvature occurs one derivative later.

In fact let \(m:\Lambda^+\otimes S^+\to S^+\) be Clifford contraction
and \(\Pi=1-m^*m/6\) its kernel projection. With an orthonormal
self-dual frame \(\omega_a\), commuting spin derivatives and subtracting
the Clifford trace gives
\[
 \Pi(d^\nabla P_{\rm tw}\psi)^+
 =-\tfrac12\sum_a\omega_a\otimes c(W^+\omega_a)\psi,\qquad
 |\Pi(d^\nabla P_{\rm tw}\psi)^+|^2
 =\tfrac12 |W^+|_{\rm HS}^2|\psi|^2.
\]
The scalar-curvature and \(D^2\psi\) terms belong to \(\mathrm{im}\,m^*\);
symmetry and trace freeness put the displayed Weyl term in \(\ker m\).
The second equality uses \(|c(\omega)\psi|^2=2|\omega|^2|\psi|^2\)
for real self-dual \(\omega\) \cite{Baker2026DiracYamabe}.
Wherever \(\psi\ne0\), dividing the squared defect by \(|\psi|^2\)
recovers exactly half the chiral density. An independent bound on that
quotient would give curvature control, whereas a small first-order
spinor Rayleigh quotient alone permits concentration.

\section{Transporting the estimates to a degenerating source}
\label{sec:marked-accounting}

Let \(g_i\) be smooth unit-volume positive Yamabe minimizers on one
fixed, oriented smooth homotopy four-sphere \(M\). Write
\(\mu_i=|W^+_{g_i}|_{\rm HS}^2dV_{g_i}\) and pass to a subsequence
with \(\mu_i(M)\to e\in[T,2T-\delta]\), \(T=4\pi^2,\ \delta>0\).
For finitely many models indexed by \(j\), suppose disjointly marked
compact truncations \(K^j_L\) on the source have masses tending, at
each fixed \(L\), to \(\mu^j(K^j_L)\). Suppose these compact sets
exhaust their models and put \(e_j=\mu^j(\bigcup_L K^j_L)\).
Writing \(R_{i,L}\) for the actual uncovered region, subtraction gives
\begin{equation}\label{eq:finite-remainder}
 \lim_{L\to\infty}\lim_{i\to\infty}\mu_i(R_{i,L})=e-\sum_j e_j .
\end{equation}
The inner limit uses finitely many mass limits at a fixed truncation.
For example, in oriented orthonormal frames on a fixed compact
let \(w_{ab,i}\) be the nine components of the self-dual Weyl block
of the pulled-back metric and \(J_i\) its volume density relative to
a fixed positive background. The actual integral has integrand
\(\sum_{a,b}|w_{ab,i}|^2 J_i\). Convergence of the metric two-jets
with uniform ellipticity on the compact gives convergence of these
components in convergent frames and of the densities. Boundedness
then gives the required mass limit by integration. A weaker comparison
must supply its own convergence argument. On a subdivided collar, a
sufficient criterion bounds the total residual by the sum of the
nonnegative errors on its pieces and the mass missed by them, and
takes that upper bound to zero in the same order. Smallness on each
of a growing number of pieces need not imply that estimate
\cite{Baker2026MarkedNecks,Baker2026DiracTransmission}.

One smooth cut can already suffice without a decomposition of all the
energy. Suppose a whole marked core has limiting uncapped energy \(x>T\) and its
boundary compares through buffers on both sides to an ordinary
coordinate sphere of a smooth compactification. Assume the caps
preserve positive Yamabe sign and the exterior cap and splice cost
is at most \(\eta\), with the capped exterior energy in the limit superior at most \(e-x+\eta\).
For \(\eta\le\delta\) this bound is strictly below \(T\).
The capped exterior is a homotopy sphere when the embedded \(S^3\)
separates a homotopy sphere: van Kampen, Mayer--Vietoris and capping
give simple connectivity and sphere homology. The strict sphere theorem
recognizes it as standard for all sufficiently late cuts. Transport
the exterior's parametrized cap into the standard sphere. The smooth
disk theorem identifies it with the prescribed standard hemisphere
with its parametrization; collar adjustment gives the corresponding
identification of complements. Gluing the latter along the original
mark therefore recovers the smooth type of the capped core
\cite[Theorem~C]{Palais1960Extending}\cite{Baker2026MarkedNecks}. If \(e=I^+(M)\), the infimum over positive
Yamabe classes on this smooth \(M\), any capped core with energy
strictly below \(e\) is now impossible. For an almost critical sequence
\(\int_M|B_{g_i}|\,dV_{g_i}\to0\), a complete smooth nonflat ALE
blow-up at scales tending to zero, with whole oriented smooth
truncation comparisons and the differentiated end decay required for
smooth compactification, has such a strict energy loss
\cite{Baker2026MarkedNecks}. A closed nonround conformal base can instead
carry all \(e\).

For the latter conclusion, take a smooth closed oriented homotopy
sphere with finitely many ordinary punctures, and supply compatible
orientation-preserving smooth embeddings of whole truncations. Write \(\Phi_i^*g_i=u_i^2m_i\);
require \(m_i\to h\) locally in \(C^0\cap W^{1,4}\) strongly, a
locally bounded \(L^2\) scalar curvature of \(m_i\), a smooth filled
metric \(h\), and \(\int_M|B_{g_i}|\,dV_{g_i}\to0\).
When a fixed compact sees positive limiting chiral energy, the
variational comparison gives a nonround Bach-flat base with energy
above \(T\). The cap argument identifies its smooth type. At
\(e=I^+(M)\), lower semicontinuity and the infimum force equality of
energies. The physical volume then fills the base, and the physical
metrics converge after subsequence in measured Gromov--Hausdorff distance to a smooth
volume-one Yamabe representative. The curvature energy has no residual
on exhaustion. Each conclusion concerns the specified markings and
metrics \cite{Baker2026MarkedLimits}.

\section{Outstanding geometric questions}

There is a source of markings before a whole base has been found.
For every pair \(\epsilon,\zeta>0\), after a subsequence one can mark
finitely many disjoint smooth standard four-balls in the original
\(M\) with total physical volume at least \(1-\epsilon\) and individual
physical full-curvature integrals below \(\zeta\). The construction
normalizes a positive-Ricci representative by the half-mass radius of
the physical measure. It transfers regular balls in that distance to
the smooth sources \cite{Baker2026MarkedLimits}. The open step is to
control the first derivatives in compatible markings and to locate
the chiral measure relative to their boundaries. It may concentrate
in a set of small physical volume. A usable separating cut needs
quantitative control on both sides, in the gauges used for sewing.

For the spinor approach the remaining estimate concerns the
normalized differentiated defect on enough of the actual source,
including control of uncovered chiral mass. For a closed nonround
critical minimizer even full recovery leaves its classification open.
Finally the entire subcritical argument requires a positive Yamabe
class with \(E_+<8\pi^2\) on the smooth manifold under investigation.
Producing such an entry class on an arbitrary smooth homotopy
four-sphere is a separate existence problem.

\bibliographystyle{unsrturl}
\bibliography{references}
\end{document}
